% Projective Temporal Ordering and Cumulative Projected Capacity
% Short note — O-series companion
% Status: draft — all sections written

\documentclass[11pt,a4paper]{article}
\usepackage{amsmath,amssymb,amsthm}
\usepackage{hyperref}
\usepackage{booktabs}
\usepackage{graphicx}

\newtheorem{definition}{Definition}
\newtheorem{proposition}{Proposition}
\newtheorem{remark}{Remark}

\title{Projective Temporal Ordering and Cumulative Projected Capacity:\\
Numerical Evidence from the Spectral Admissibility Cascade}
\author{J\'er\^ome Beau}
\date{2026}

\begin{document}
  \maketitle

% ------------------------------------------------------------------
  \begin{abstract}
    We identify a monotone proxy for temporal ordering within the
    Cosmochrony spectral admissibility programme.
    The natural observable along the BFS cascade is not the residual
    capacity $\sigma_{\mathrm{pair}}(n)$, which decays by construction,
    but the cumulative projected capacity
    $I(n) = \sigma_{\mathrm{pair}}(0) - \sigma_{\mathrm{pair}}(n)$,
    which measures the spectral structure already projected at depth~$n$.
    For $q \in \{29, 61, 101, 151\}$, $I(n)$ is found to be strictly
    non-decreasing at every BFS step, for every conjugate pair
    $(c, q{-}c)$, with zero physically meaningful regressions observed
    across 169~pairs and 310~steps after burn-in.
    This supports the identification of temporal ordering with the monotone
    accumulation of cumulative projected capacity along the admissible
    cascade.
    The saturation plateau of $I(n)$ is interpreted as Born--Infeld
    saturation at the pair level.
  \end{abstract}

% ------------------------------------------------------------------
  \section{Motivation}

    The Cosmochrony framework identifies observable dynamics not with an
    evolution of the relational substrate~$\chi$---which is static---but
    with an iterated reprojection sequence
    $U_{t+1} = \Pi \circ \sigma(U_t)$,
    where $\Pi$ is a fixed non-injective projection and
    $\sigma(U_t) \in \Pi^{-1}(U_t)$ selects an admissible configuration
    within the fiber compatible with the current effective
    state~$U_t$~\cite{Beau2026a}.
    This formulation raises a concrete question: what scalar observable,
    defined along the cascade, plays the role of a temporal ordering
    parameter?
    A candidate must be monotone along every admissible trajectory, stable
    against fluctuations, and computable from the pipeline.
    The spectral admissibility programme produces, for each prime~$q$ and
    each conjugate pair $(c, q{-}c)$, a family of traces
    $\sigma_{\mathrm{pair}}(n)$ indexed by BFS depth~$n$~\cite{Beau2026a29}.
    We show that the correct variable is not $\sigma_{\mathrm{pair}}(n)$
    itself, but its complement
    $I(n) = \sigma_{\mathrm{pair}}(0) - \sigma_{\mathrm{pair}}(n)$.

% ------------------------------------------------------------------
  \section{The Two Variables and Why $\sigma_{\mathrm{pair}}$ Is Not the
  Right One}

    The observable $\sigma_{\mathrm{pair}}(n) = \sigma_c(n)\,\sigma_{q-c}(n)$
    measures the residual projective capacity of the conjugate pair at
    depth~$n$: it counts, normalised by~$q$, the fraction of Weil
    fingerprint vectors that remain linearly independent of the
    Gram--Schmidt span built over the preceding shells.
    By construction, $\sigma_{\mathrm{pair}}(n)$ decays as~$n$ grows and
    cannot serve as a temporal proxy, since a decreasing quantity does not
    define an ordering compatible with the accumulation of structure.
    The correct variable is its complement.

    \begin{definition}[Cumulative projected capacity]
Let $\sigma_{\mathrm{pair}}(n)$ denote the canonical pair
observable~\cite{Beau2026a29} at BFS depth~$n$.
The \emph{cumulative projected capacity} at depth~$n$ is
\begin{equation}
  \label{eq:I}
  I(n) \;:=\; \sigma_{\mathrm{pair}}(0) - \sigma_{\mathrm{pair}}(n).
\end{equation}
    \end{definition}

    $I(n)$ measures the spectral structure that has already been stably
    projected up to depth~$n$: it starts at zero and accumulates as the
    cascade progresses, with
    $I(0) = 0$ and $I(n) \to \sigma_{\mathrm{pair}}(0)$ as
    $\sigma_{\mathrm{pair}}(n) \to 0$.

    \begin{remark}
      $I(n)$ is non-decreasing if and only if $\sigma_{\mathrm{pair}}(n)$ is
      non-increasing.
      While this is expected under a dissipative interpretation of the BFS
      cascade, the present work establishes it empirically without relying on
      a formal dissipation theorem: the monotonicity of $I(n)$ is observed
      directly, step by step and pair by pair, across all tested primes.
    \end{remark}

% ------------------------------------------------------------------
  \section{Temporal Ordering Conjecture}

    \begin{proposition}[Projective temporal ordering --- numerical form]
      \label{prop:order}
      Define an ordering on effective states along the BFS cascade by
      \begin{equation}
        \label{eq:order}
        U_t \prec U_{t+1}
        \quad \Longleftrightarrow \quad
        I(U_t) \leq I(U_{t+1}).
      \end{equation}
      For $q \in \{29, 61, 101, 151\}$, this ordering is strictly observed
      along the tested BFS trajectories: $I(n)$ is strictly non-decreasing
      at every step, for every conjugate pair $(c, q{-}c)$, with no
      physically meaningful exception.
    \end{proposition}

    \begin{remark}[Status]
This is numerical evidence, not a proof.
The result establishes the ordering empirically along the tested
trajectories and does not exclude the possibility of exceptions at
larger~$q$ or along non-BFS paths.
The one-way activation property---that modes, once stably projected,
do not deactivate---remains to be established analytically.
    \end{remark}

% ------------------------------------------------------------------
  \section{Protocol and Data}

    Computations use the \texttt{o25\_paired\_pipeline.py} pipeline with
    parameters \texttt{M=50} independent block samples per pair and
    \texttt{bfs\_frac=0.99} (full graph exploration).
    Monotonicity of $I(n)$ is tested by the
    \texttt{o\_stability\_diagnostic.py} script, which verifies,
    for each prime and each pair, that no step $I(n{+}1) < I(n)$ occurs
    after the burn-in depth $n_{\mathrm{burn}} = 2$.

    \begin{table}[h]
      \centering
      \caption{Monotonicity and stability of $I(n)$ per prime.
      ``Steps'' counts BFS depths after burn-in.
      ``Mean mono'' reports whether the pair-averaged $I(n)$ is strictly
      non-decreasing.
      ``Regressions'' counts pairs with at least one physically meaningful
      decrease in $I$ (threshold: relative decrease $> 10^{-6}$ of the
      plateau value).}
      \label{tab:stability}
      \begin{tabular}{rrrrrr}
        \toprule
        $q$ & pairs & steps & mean mono & regressions & mono frac \\
        \midrule
        29 &  14 &  25 & YES & 0 & 1.0000 \\
        61 &  30 &  55 & YES & 0 & 1.0000 \\
        101 &  50 &  92 & YES & 0 & 1.0000 \\
        151 &  75 & 138 & YES & $0^*$ & $1.0000^*$ \\
        \bottomrule
      \end{tabular}
      \smallskip

      {\footnotesize
        $^*$At $q = 151$, 70 pairs exhibit floating-point fluctuations of
      order $10^{-10}$--$10^{-13}$ on the saturation plateau, where
        $I/I_{\max} = 1.000000$.
      These are ${\sim}6$ orders of magnitude below the mean positive step
      size of $4 \times 10^{-2}$ and are not counted as physical regressions.
      }
    \end{table}

    \begin{figure}[h]
      \centering
% \includegraphics[width=\linewidth]{figures/traces_q101.pdf}
      \caption{Left: residual capacity $\sigma_{\mathrm{pair}}(n)$ (strictly
      decreasing) for $q = 101$, all 50 pairs shown.
      Right: cumulative projected capacity
        $I(n) = \sigma_{\mathrm{pair}}(0) - \sigma_{\mathrm{pair}}(n)$
        (strictly non-decreasing) for the same pairs.
        The two panels are exact complements and together visualize the
        one-way activation property.
        The plateau in the right panel corresponds to Born--Infeld saturation
        at the pair level.}
      \label{fig:traces}
    \end{figure}

% ------------------------------------------------------------------
  \section{Interpretation}

    \paragraph{One-way activation.}
      The zero-regression result means the admissible BFS cascade behaves as
      a one-way activation process in spectral space.
      Once a mode contributes stably to the Gram--Schmidt span, it does not
      leave: the spectral structure accumulated at depth~$n$ is preserved at
      all subsequent depths.
      This is the empirical content of the stability condition required for
      the identification of $I(n)$ as a temporal proxy: an ordering defined
      by $I$ is well-posed only if $I$ never reverses.

    \paragraph{Saturation plateau and Born--Infeld bound.}
      $I(n)$ reaches its maximum $\sigma_{\mathrm{pair}}(0)$ when the
      residual capacity is exhausted.
      This plateau is interpreted as Born--Infeld saturation at the pair
      level: once the projective flux budget of the conjugate pair is fully
      consumed, no additional spectral structure can be resolved within the
      effective description.
      Saturation is irreversible in all tested cases; this is consistent with
      the bounded-flux constraint of the Cosmochrony
      framework~\cite{Beau2026a} but is not derived from it here.

    \paragraph{Candidate arrow of time.}
      The strict monotonicity of $I(n)$ provides a candidate ordering
      compatible with an arrow of time, without invoking thermodynamic entropy
      or a background temporal coordinate.
      The ordering is projective: it reflects the irreversible accumulation
      of stably activated spectral structure along the admissible cascade.
      Whether this ordering coincides with the physical time parameter of the
      full framework~\cite{Beau2026a} is left as an open question; the present
      result establishes only its monotone and stable character along the
      tested BFS trajectories for $q \in \{29, 61, 101, 151\}$.

% ------------------------------------------------------------------
  \section{Open Questions}

  \begin{enumerate}

    \item Does strict monotonicity of $I(n)$ hold for all admissible
    primes~$q$, or are there exceptions at large~$q$?
    Extension to $q \in \{211, 307, 401\}$ is pending; preliminary data
    at $q = 151$ are consistent with the property
    (Table~\ref{tab:stability}).

    \item Can the one-way activation property---$\sigma_{\mathrm{pair}}(n)$
    non-increasing along the full cascade---be proved analytically from
    the structure of the BFS expansion and the admissibility constraints
    of~$\chi$?

    \item What is the precise relationship between $I(n)$ and the capacity
    exponent $\delta_{\mathrm{pair}}$~\cite{Beau2026a29}?
    Is $I(n)$ sensitive to the value of $\delta_{\mathrm{pair}}$ or only
    to its sign?

    \item Does the saturation depth (the~$n$ at which $I(n)$ plateaus)
    carry physical meaning independent of~$q$?

    \item What is the precise relationship between BFS shell depth and
    effective resolution-cell entry?
    Preliminary numerical evidence suggests that the characteristic
    shell~$n^*$ at which $\sigma_{\mathrm{pair}}(n)$ first undergoes a
    qualitative knee increases with~$q$ over the tested cases
    $q \in \{29, 61, 101\}$, with approximate values
    $n^*(29) \approx 4$, $n^*(61) \approx 6$, $n^*(101) \approx 8$.
    This suggests that BFS shells do not define resolution cells, but
    may reveal the depth at which such cells become effective.
    The structural status, scaling law, and analytical derivation of
    this relation are left open.

  \end{enumerate}

% ------------------------------------------------------------------
  \section{Conclusion}

  The observable $\sigma_{\mathrm{pair}}(n)$ measures residual projective
  capacity and decays along the BFS cascade; its complement
  $I(n) = \sigma_{\mathrm{pair}}(0) - \sigma_{\mathrm{pair}}(n)$
  measures cumulative projected capacity and grows.
  Across $q \in \{29, 61, 101, 151\}$, covering 169~pairs and
  310~post-burn-in steps, $I(n)$ is strictly non-decreasing at every BFS
  step and for every conjugate pair, with no physically meaningful
  regression.
  This establishes $I(n)$ as a robust monotone proxy for temporal ordering
  along the admissible spectral cascade, consistent with the identification
  of time with the irreversible accumulation of stably projected structure
  in the Cosmochrony framework~\cite{Beau2026a}.
  The one-way activation property observed here is a numerical result;
  its analytical derivation from the admissibility constraints of~$\Pi$
  and the scaling of $n_1(q)/q$~\cite{Beau2026a29} remains the central
  open direction.

% ------------------------------------------------------------------
  \bibliographystyle{plain}
  \bibliography{cosmochrony-bibliography}

\end{document}
